% !PRINTED 661
\weberpage{661}
%<<CONT
in which $\xi$ runs through the sequence of numbers determined more closely in \S.\,172.
Let $\mathfrak{p}_\xi^{k_\xi}$ be the power of $\mathfrak{p}_\xi$ contained in $\varphi$,
so that, if we set
\[
\varphi = \chi \prod^{\xi} \mathfrak{p}_\xi^{k_\xi}
\]
then $\chi$ is coprime to $p$ and is related to no primes other than
those to which $\varphi$ itself is related (see the conclusion of
\S.\,173). From this it follows that
\weq{5}{\varphi_n = \chi_n \prod^{\xi} \mathfrak{p}_{n\xi}^{k_\xi},}
and here we again let $n$ run through the sequence of numbers $\xi$.
If we further set
\[
\psi = \prod^{\xi} \chi_\xi,
\]
then $\psi$ too is coprime to $p$ and is related to no primes other
than those to which $\varphi$ is related, and from (5), using (4), we obtain:
\weq{6}{\prod^{\xi} \varphi_\xi = \psi\, p^{\Sigma k_\xi},}
and from this it follows that the functional $\psi$ satisfies the same two
conditions (1), (2) as $\varphi$.

It appears further that, if the $\varphi_\xi$ lie in the principal class,
then $\psi$ too is contained in the principal class.

If, however, $p-1$ is not divisible by $q$, then we can also
infer the converse.

For by (2) $\varphi_\xi$ is equivalent to $\varphi^\xi$, and consequently:
\[
\begin{gathered}
\prod^{\xi} \varphi_\xi \ \text{equivalent}\ \varphi^{\Sigma\xi}, \\
\text{equivalent}\ \psi,
\end{gathered}
\]
and if $\psi$ is associated with a number, the same holds
of $\varphi^{\Sigma\xi}$. But by the theorem of \S.\,172, if $p-1$ is not
divisible by $q$, then $\Sigma \xi$ is also not divisible by $q$, and consequently
 $\varphi$ itself is likewise contained in the principal class.

This consideration we can now repeat, if $\psi$ is related to
still further primes that are not congruent to 1 with respect to the
 modulus $q$, and we thus arrive at the
theorem:

\gesperrt{Theorem a) is proved in general if it can be
proved under the assumption
that $\varphi$ is related only to such primes as
are congruent to 1 with respect to the modulus $q$.}
